% Generated draft: requires mathematical and source review.
\documentclass[11pt]{article}
\usepackage[margin=25mm]{geometry}
\usepackage{fontspec}
\setmainfont{Latin Modern Roman}
\usepackage{amsmath,amssymb,mathtools,unicode-math}
\setmathfont{Latin Modern Math}
\usepackage{enumitem,needspace,xurl,hyperref}
\setlength{\parindent}{0pt}
\setlength{\parskip}{5pt}
\setlength{\emergencystretch}{3em}
\newcommand{\N}{\mathbb{N}}
\newcommand{\Z}{\mathbb{Z}}
\newcommand{\Q}{\mathbb{Q}}
\newcommand{\R}{\mathbb{R}}
\newcommand{\C}{\mathbb{C}}
\newcommand{\F}{\mathbb{F}}
\newcommand{\A}{\mathbb{A}}
\newcommand{\PP}{\mathbb{P}}
\newcommand{\E}{\mathbb{E}}
\newcommand{\eps}{\varepsilon}
\newcommand{\dd}{\,\mathrm{d}}
\begin{document}
\section*{UnsolvedMath ID: 800606}
\textbf{Source problem code:} catalogue entry
\subsection{Definitions and mathematical statement}
For a topological space \(X\), a property \(P\) is \emph{hereditary} if every subspace of every space having \(P\) also has \(P\). More explicitly, relative to an intended class \(\mathcal C\) of spaces, write
\[
H_{\mathcal C}(P):\quad
\forall X\in\mathcal C\;\forall Y\subseteq X,
\qquad P(X)\Longrightarrow P(Y),
\]
where \(Y\) carries the subspace topology.

The four properties named in the supplied record are
\[
\begin{aligned}
P_1&=\mathsf{Split}(\mathrm{B}_{\Omega},\mathrm{B}_{\Lambda}),\\
P_2&=\mathsf{Split}(\mathrm{B}_{\Omega},\mathrm{B}_{\Omega}),\\
P_3&=\mathsf{Split}(\mathrm{B}_{\mathrm{T}},\mathrm{B}_{\mathrm{T}}),\\
P_4&=\mathsf{Split}(\mathrm{B}_{\mathrm{T}},\mathrm{B}_{\Gamma}).
\end{aligned}
\]
The question is whether
\[
\bigvee_{i=1}^{4}H_{\mathcal C}(P_i)
\]
holds: is at least one of these four properties hereditary?

The supplied record does not specify \(\mathcal C\) or define the cover classes and splitting operator. The following supplies a conditional interpretation using customary cover terminology; these conventions require confirmation against the cited source. For a space \(X\), consider countable families \(\mathcal U\) of Borel subsets of \(X\) such that \(\bigcup\mathcal U=X\) and \(X\notin\mathcal U\). Such a cover is \emph{large} if every point belongs to infinitely many members; it is an \(\omega\)-cover if every finite subset of \(X\) is contained in some member; and it is a \(\gamma\)-cover if \(\mathcal U\) is infinite and each point belongs to all but finitely many members. A \(\tau\)-cover is a large cover satisfying, for every \(x,y\in X\),
\[
\bigl|\{U\in\mathcal U:x\in U,\ y\notin U\}\bigr|<\infty
\quad\text{or}\quad
\bigl|\{U\in\mathcal U:y\in U,\ x\notin U\}\bigr|<\infty.
\]
Under this interpretation, \(\mathrm{B}_{\Lambda}\), \(\mathrm{B}_{\Omega}\), \(\mathrm{B}_{\Gamma}\), and \(\mathrm{B}_{\mathrm{T}}\) denote, respectively, these Borel large, \(\omega\)-, \(\gamma\)-, and \(\tau\)-covers.

One possible precise splitting convention is that \(X\) satisfies \(\mathsf{Split}(\mathcal A,\mathcal B)\) when
\[
\forall\mathcal U\in\mathcal A(X)\;\exists\mathcal V_0,\mathcal V_1\subseteq\mathcal U,
\qquad
\mathcal V_0\cap\mathcal V_1=\varnothing,
\quad
\mathcal V_0,\mathcal V_1\in\mathcal B(X).
\]
A partition convention additionally requires \(\mathcal V_0\cup\mathcal V_1=\mathcal U\). The supplied record does not resolve which convention is intended. Thus the displayed heredity question preserves the source question, while its fully instantiated topological formulation remains conditional on the ambient class and these definitions. No mathematical resolution is asserted.
\subsection{Short English statement}
Is at least one of the four listed Borel cover splitting properties inherited by every subspace? The intended class of spaces and exact cover and splitting conventions need confirmation from the cited source.
\subsection{Sources}
\begin{itemize}
\item Local problem registry: permanent UnsolvedMath ID \(800606\), source problem code \texttt{catalogue entry}, AMR Open Problem Lists collection. The supplied registry attributes the problem to Tsaban, with proposed year \(2006\), and labels it \emph{open}; that status is supplied metadata, not an independently verified assessment.
\item Supplied source: Tsaban, \emph{Selection Principles and Special Sets of Reals: Open Problems} (2006), Problem 34, \url{https://arxiv.org/abs/math/0609601}. The record reports access on 2026-10-16 and extraction by \texttt{arxiv\_latex}. These provenance claims are supplied and have not been independently verified here.
\item The supplied title and background include the reference/context label \texttt{[split]}. No bibliographic expansion of that label or additional literature reference is supplied.
\end{itemize}
\end{document}
